% !TEX root = ../../main.tex
\chapter*{Acknowledgements}
\addcontentsline{toc}{chapter}{Acknowledgements}

These notes grow through the shared work of students studying mathematical
analysis. We are grateful to everyone who contributes explanations,
checks proofs, suggests corrections, or improves the presentation.

We thank Professor Yisong Yang of NYU Courant for his teaching and the
mathematical framework on which these notes are based.
% Name and course affiliation verified against NYU Courant's Fall 2026 listing.
% https://math.nyu.edu/dynamic/courses/graduate-course-descriptions/tandon/fall-2026/

\section*{Student contributors}

\begin{itemize}
  \item Lucas Barbosa --- handwritten notes and course structure
    for Mathematical Analysis I.
\end{itemize}

\section*{General acknowledgements}

OpenAI Codex assisted with the Makefile support for selecting individual
homework handouts, and with transcription, mathematical editing, and TikZ
diagrams for the foundations chapter and the shared homework exercises.
It also assisted with importing Homework 03, organizing its shared
statements, and completing the solutions and diagrams for Questions 2--8.
It also assisted with adapting the formula-sheet layout and checking its
sequence criteria and summation identities.
It assisted with transcribing the sequence, series, and power-series
notes, checking their estimates and endpoint conditions, and drawing
supporting diagrams.
It also assisted with organising the formula sheet in the book's topic order
and providing a standalone entry that shares its content, auditing coverage
of Homework 01--04, and collecting reusable proof techniques.
It also assisted with configuring and auditing alphabetic bibliography
formatting and checking citations in the book and standalone homework documents.
It also assisted with updating the book's license notices and official
Creative Commons badge.
It also assisted with transcribing the differentiation notes,
reviewing their local arguments, and redrawing the supplied geometric sketches.
It assisted with the first-principles derivative diagram and its
manuscript-based historical margin note.
It also assisted with importing the Homework 06 statements, preserving their
subparts and hints, and checking their shared placement and standalone handout.
It assisted with the accompanying function-space and convergence background,
series products, and the construction of exponential and logarithmic functions.
It assisted with reconciling additional handwritten pages, transcribing the
moving-centre Taylor argument, convexity and its inequalities, analytic
functions, and Bernstein's polynomial approximation proof. It also prepared
supporting diagrams and the shared Homework 07 statements and handout.
It also assisted with importing Homework 05, checking its topic coverage
and prerequisites, and maintaining the shared exercise placement map.
At the author's request, it completed Homework 05 Questions 1 and 3--10
and their supporting diagrams for review.
It also assisted with importing and placing the shared Homework 04
statements, transcribing and reviewing the author's solutions, and preparing
their standalone handout and appendix entries.
It also assisted with the number-systems overview and its nested-box diagram.
It assisted with transcribing the author's interval-halving proof of the
continuum cardinality and checking its binary and ternary encodings.
It also assisted with the comparison of touching and separated interval codes
and the historical notes on the middle-thirds construction.
It also assisted with the historical context and names of the
Cantor--Bernstein theorem.
It assisted with proofs of that theorem, the general arithmetic--geometric
mean inequality, the function-limit laws, Cauchy's mean value theorem, and
Taylor's theorem with Lagrange remainder.
It also assisted with illustrating the two reciprocal chains that make room
for the endpoints of an interval.
It also assisted with placing the interval-cardinality exercises beside
the bijections they reinforce.
It also assisted with the function-cardinality comparison and its
set-and-arrow diagrams.
It also assisted with the decimal-approximation proof of rational density
and the successive zooms around the square root of two.
It assisted with aligning solution links to exercise headings and checking
their placement throughout the book.
It also assisted with transcribing the continuity and function-sequence
notes, completing abbreviated arguments for review, and drawing their
supporting diagrams.
It also assisted with clarifying the rational-to-real extension for
continuous additive functions.
It also assisted with stating and proving Euclid's lemma and connecting
its later divisibility applications to that statement, including margin
notes on its history and role in unique prime factorisation.
It also assisted with the assumed-knowledge recall of divisibility,
greatest common divisors, and lowest-terms contradictions.
It also assisted with the inventory of historically attributable material,
historical margin notes, portrait reproduction credits, and the AMS
bibliography configuration.
It also assisted with moving solution navigation to italic links at the
ends of exercises and checking their appendix destinations.
It also assisted with configuring GitHub Actions to publish versioned PDFs
of the main book.
It assisted with a book-wide editorial review, the mathematical notation
reference and introduction, clearer proof explanations, and removal of
unrelated typesetting specimens.
It also assisted with transcribing the author's supplied differentiation
solutions, reviewing their mathematical steps through focused amendments,
and preparing supporting diagrams while preserving the author's arguments.
